\section{Reducing Code Conversion to Codespace Projection}
\label{sec:reducing_code_conversion_to_codespace_projection}
We first give two definitions.
\emph{Code conversion} is a quantum circuit,
which may include measurement and feedforward,
that changes encoding and preserves every logical state.
This is more general than code switching
as it may involve conversion into the same code family
e.g.\ surface code to surface code.
Though MSC uses a code conversion technique called grafting in its escape stage,
later MSC variants demonstrate that
other code conversion techniques can be used:
	measurement-based growth is used in~\cite{Sahay2026,Claes2025,Chen2026},
	teleportation via lattice surgery is used in~\cite{Itogawa2025,Hirano2025a,Thorvaldson2026},
	and local unitary transformations are used in~\cite{Itogawa2025,Sahay2026}.

Associated to each detector is a \emph{detecting region} which
(for stabiliser circuits)
is a set of circuit spacetime locations
each labelled by X, Y, or Z~\cite{McEwen2023},
and shows how the Pauli measured by that detector
evolves through the circuit at each timeslice.

\begin{lem}
\label{lem:insert_projector}
Let $V$ be a code conversion circuit
from an $\llbracket n,k\rrbracket$ input code generated by $\mathcal G$.
Suppose $n -k$ detecting regions have
independent timeslices $P_j$ at the input of $V$,
with eigenvalue $+1$ when the corresponding detector is unflipped.
Then in the absence of noise,
for an arbitrary mixed state $\rho$ entering $V$,
the probability those detectors
and any others within or after $V$ are unflipped is
\begin{equation}
	\pr(\v s =\v 0)
	=\tr(\upPi_{\mathcal G} \rho).
	\label{eq:insert_projector}
\end{equation}
Conditioned on this event,
inserting the codespace projector $\upPi_{\mathcal G}$
immediately before $V$ leaves the output state unchanged.
\end{lem}
\begin{proof}
In the absence of noise,
each detector within or after $V$ can be pulled back along its detecting region
to its timeslice $Q$ at the input of $V$.
Code conversion preserves every logical state,
so $Q \in `<\mathcal G>$.
The $n -k =\rk `<\mathcal G>$ timeslices $P_j$ in particular are independent,
so $`<P_j> =`<\mathcal G>$.
Conditioning on their corresponding $n -k$ detectors being unflipped
projects $\rho$ onto the simultaneous $+1$ eigenspace of $\mathcal G$,
making redundant the condition any other detector is unflipped.
This proves \cref{eq:insert_projector}.

Conditioned on this event:
the projection already occurs within or after $V$,
so inserting it beforehand leaves the output unchanged.
\end{proof}